% =============================================================================
% 纸笔外语学习 - LaTeX 环境定义
% Paper-and-Pencil Language Learning - LaTeX Environment Definitions
% =============================================================================

% 基础宏包依赖
\usepackage{framed}
\usepackage{xcolor}
\usepackage{enumitem}
\usepackage[normalem]{ulem}
\usepackage{array}

% =============================================================================
% 通用框架
% =============================================================================

% 任务框环境（采用无框横幅卡片，彻底消除与 multicols 输出例程冲突，允许内部习题跨栏跨页无损分页）
\newenvironment{taskbox}[2][]
  {%
    \par\vspace{0.6em}\noindent
    \begingroup
    \small
    \noindent\colorbox{blue!7}{%
      \parbox{\dimexpr\linewidth-2\fboxsep\relax}{%
        \bfseries\color{blue!80!black}#2%
        \if\relax\detokenize{#1}\relax\else\par\vspace{0.2em}\normalfont\small\textit{#1}\fi
      }%
    }%
    \par\nopagebreak\vspace{0.4em}%
    \endgroup
  }%
  {\par\vspace{0.6em}}

% 侧边提示框
\newenvironment{sidebox}[1]
  {%
    \par\vspace{0.4em}\noindent
    \begin{minipage}{\linewidth}%
      \small\textbf{#1}\par\vspace{0.2em}%
  }%
  {\end{minipage}\par\vspace{0.4em}}

% 答案空间命令（单源双输出：教师版压缩，学生版留白，增加弹性伸缩解决硬性断页）
\ifdefined\TeacherVersion
  \NewDocumentCommand{\ansspace}{o g}{\par\vspace{0.4ex}}
\else
  \NewDocumentCommand{\ansspace}{o g}{%
    \IfValueTF{#1}{%
      \par\vspace{#1 plus 0.2\dimexpr#1\relax minus 0.3\dimexpr#1\relax}%
    }{%
      \IfValueTF{#2}{%
        \par\vspace{#2 plus 0.2\dimexpr#2\relax minus 0.3\dimexpr#2\relax}%
      }{%
        \par\vspace{2.5cm plus 0.5cm minus 0.8cm}%
      }%
    }%
  }
\fi
\newcommand{\ansspaceend}{}

% 填空宏（教师版使用 uline 允许长句答案自然断行，避免硬性定宽造成行溢出；学生版呈现空白下划线）
\ifdefined\TeacherVersion
  \newcommand{\fillin}[2][]{%
    \if\relax\detokenize{#1}\relax
      \uline{\hspace{#2}}%
    \else
      \textcolor{blue!85!black}{\uline{\hspace{0.25em}\bfseries#1\hspace{0.25em}}}%
    \fi
  }
\else
  \newcommand{\fillin}[2][]{%
    \uline{\hspace{#2}}%
  }
\fi

% =============================================================================
% 阅读理解任务
% =============================================================================

\newenvironment{readingtask}[2][]{%
  \begin{taskbox}[#1]{【阅读理解】#2}%
  \textbf{阅读以下文本，然后回答问题：}%
  \vspace{0.8em}%
  \begin{readingtext}%
}{%
  \end{readingtext}%
  \end{taskbox}%
}

% 阅读文本容器
\newenvironment{readingtext}{}{}

% 问题列表
\newenvironment{questions}{\begin{enumerate}[label=\arabic*, leftmargin=*]}{\end{enumerate}}

% 问题编号命令（带分值）
\newcommand{\question}[2][1]{\item [\textbf{(#1 分)}] #2}

% =============================================================================
% 语法操练
% =============================================================================

\newenvironment{grammardrill}[2][]{%
  \begin{taskbox}[#1]{【语法要点】#2}%
  \if\relax\detokenize{#1}\relax\else\textbf{规则：}#1\par\vspace{0.5em}\fi%
}{%
  \end{taskbox}%
}

% 对比说明
\newcommand{\contrast}[1]{%
  \vspace{0.3em}%
  \textbf{对比：}#1%
  \vspace{0.3em}%
}

% 练习集合
\newenvironment{exercises}{\begin{enumerate}[label=\arabic*, leftmargin=*]}{\end{enumerate}}

% =============================================================================
% 写作支架
% =============================================================================

\newenvironment{writingscaffold}[2][]{%
  \begin{taskbox}[#1]{【写作任务】#2}%
  \textbf{任务：}#1\par\vspace{0.5em}%
}{%
  \end{taskbox}%
}

% 场景描述
\newcommand{\scenario}[1]{%
  \textbf{情景：}#1%
  \vspace{0.3em}%
}

% 任务目标
\newcommand{\taskgoal}[1]{%
  \textbf{目标：}#1%
  \vspace{0.3em}%
}

% 结构指南
\newenvironment{structureguide}{%
  \textbf{结构建议：}%
  \begin{itemize}[leftmargin=*, nosep]%
}{%
  \end{itemize}\vspace{0.3em}%
}
\let\structureguidestart\structureguide
\let\structureguideend\endstructureguide

% 有用短语
\newenvironment{usefullanguage}{%
  \textbf{Useful language:}%
  \begin{itemize}[leftmargin=*, nosep]%
}{%
  \end{itemize}\vspace{0.3em}%
}
\let\usefullanguagestart\usefullanguage
\let\usefullanguageend\endusefullanguage

% 范文段落
\newenvironment{modelparagraph}{%
  \vspace{0.5em}%
  \textbf{范文示例：}%
  \begin{quote}%
}{%
  \end{quote}\vspace{0.3em}%
}
\let\modelparagraphstart\modelparagraph
\let\modelparagraphend\endmodelparagraph

% 写作时间建议
\newcommand{\writingtime}[1]{%
  \textbf{建议用时：}#1%
  \vspace{0.3em}%
}

% 检查清单
\newenvironment{checklist}{%
  \textbf{自查清单：}%
  \begin{itemize}[leftmargin=*, nosep]%
}{%
  \end{itemize}\vspace{0.3em}%
}
\let\checkliststart\checklist
\let\checklistend\endchecklist

% =============================================================================
% 体裁分析
% =============================================================================

\newenvironment{genreanalysis}[1]{%
  \begin{sidebox}{【体裁分析】#1}%
}{%
  \end{sidebox}%
}

% =============================================================================
% 翻译与转换训练（v2.0 增强版）
% =============================================================================

% 翻译任务环境
\newenvironment{translationtask}[2][]{%
  \begin{taskbox}[#1]{【汉英翻译训练】#2}%
  \textbf{Translate the following Chinese text into English.}%
  \vspace{0.8em}%
}{%
  \end{taskbox}%
}

% 中文原文框
\newenvironment{chinesetext}{%
  \begin{quote}
  \small
}{%
  \end{quote}%
}

% 关键词汇对照
\newenvironment{keywords}{%
  \textbf{关键词汇对照：}\begin{itemize}[leftmargin=*, nosep]%
}{%
  \end{itemize}\vspace{0.3em}%
}
\let\keywordsstart\keywords
\let\keywordsend\endkeywords

% 语法焦点提示
\newcommand{\syntaxfocus}[1]{%
  \vspace{0.3em}%
  \teacheronly{\textcolor{gray!70!black}{\textbf{【语法焦点】}#1}}%
}

% 字数要求
\newcommand{\wordlimit}[1]{%
  \vspace{0.3em}%
  \teacheronly{\textit{(Target: #1 words)}}%
}

% 变体答案支持
\newcommand{\alternativeanswer}[1]{%
  \teacheronly{%
    \vspace{0.3em}%
    \textbf{【替代答案】}#1%
  }%
}

% 可接受变化范围
\newcommand{\acceptablevariations}[1]{%
  \teacheronly{%
    \vspace{0.3em}%
    \textbf{【可接受变体】}#1%
  }%
}

% 评分要点
\newcommand{\scoringpoints}[1]{%
  \teacheronly{%
    \vspace{0.3em}%
    \textbf{【评分要点】}%
    \begin{itemize}[leftmargin=*, nosep]
    #1
    \end{itemize}%
  }%
}

% v2.0 旧版翻译桥梁（保留向后兼容）
\newenvironment{translationbridge}{%
  \begin{taskbox}{【汉译英练习】}%
}{%
  \end{taskbox}%
}

\newcommand{\chinesesentence}[1]{%
  \textbf{中文：}#1%
  \vspace{0.3em}%
}

\newenvironment{focuspoints}{%
  \textbf{注意：}%
  \begin{itemize}[leftmargin=*, nosep]%
}{%
  \end{itemize}\vspace{0.3em}%
}
\let\focuspointsstart\focuspoints
\let\focuspointsend\endfocuspoints

\newcommand{\translate}[1]{%
  \vspace{0.3em}%
  \textbf{#1}%
  \ansspace[6cm]%
}

\newcommand{\modelanswer}[1]{%
  \studentonly{}%
  \teacheronly{%
    \vspace{0.3em}%
    \textbf{参考译文：}#1%
  }%
}

% =============================================================================
% 偏误分析
% =============================================================================

\newenvironment{erroranalysis}[1]{%
  \begin{taskbox}{【偏误诊断】#1}%
}{%
  \end{taskbox}%
}

% 错误标注（划掉 + 正确答案提示）
\newcommand{\error}[2]{%
  \uline{\sout{#1}}%
  \studentonly{}\teacheronly{\textcolor{red!70!black}{[#2]}}%
}

\newcommand{\explain}[1]{%
  \vspace{0.3em}%
  \textit{【说明】#1}%
  \vspace{0.3em}%
}

% =============================================================================
% 交际模拟
% =============================================================================

\newenvironment{writteninteraction}[2][]{%
  \begin{taskbox}[#1]{【书面交际模拟】#2}%
  \textbf{交际场景：}#1\par\vspace{0.5em}%
}{%
  \end{taskbox}%
}

\newcommand{\participant}[1]{%
  \textbf{#1}%
}

\newcommand{\inputprovided}[1]{%
  \vspace{0.3em}%
  \textbf{提供材料：}%
  \begin{quote}%
  #1%
  \end{quote}%
  \vspace{0.3em}%
}

% =============================================================================
% 语言要点卡
% =============================================================================

% 语言要点卡（无框横幅卡片，强制末尾换段，禁止标题后孤行断页）
\newenvironment{languagecard}[1]{%
  \par\vspace{0.6em}\noindent
  \begingroup
  \small
  \noindent\colorbox{gray!12}{%
    \parbox{\dimexpr\linewidth-2\fboxsep\relax}{%
      \bfseries\color{black}【语言要点】#1%
    }%
  }%
  \par\nopagebreak\vspace{0.4em}%
  \endgroup
}{%
  \par\vspace{0.6em}%
}

\newcommand{\langrule}[1]{%
  \textbf{规则：}#1%
  \vspace{0.3em}%
}
\let\gramrule\langrule

% =============================================================================
% 标记与标签
% =============================================================================

% CEFR 等级标记
\newcommand{\cefr}[1]{\textsuperscript{\textsf{CEFR #1}}}

% 体裁标记
\newcommand{\genre}[1]{\textbf{[体裁：#1]}}

% 字数要求
\newcommand{\wordcount}[1]{\textbf{[字数：#1]}}

% 难度标记
\newcommand{\difficulty}[1]{\textcolor{blue!70!black}{\small\textbf{难度：}#1}}

% 使用短语框
\newcommand{\usefulphrase}[1]{%
  \vspace{0.3em}%
  \noindent\framebox[\linewidth][l]{%
    \parbox[t]{\dimexpr\linewidth-2\fboxsep-2\fboxrule\relax}[t]{%
      \small\textbf{Useful phrase:} #1%
    }%
  }%
  \vspace{0.3em}%
}

% =============================================================================
% 可见性控制（单源双输出：基于 environ 宏包实现完美剥离）
% =============================================================================
\usepackage{environ}

% 仅学生可见
\newcommand{\studentonly}[1]{%
  \ifdefined\TeacherVersion\else#1\fi
}

% 仅教师可见
\newcommand{\teacheronly}[1]{%
  \ifdefined\TeacherVersion#1\fi
}

\ifdefined\TeacherVersion
  % 参考答案环境（教师版醒目显示，带弹性断行）
  \NewEnviron{solution}[1][参考答案]{%
    \par\vspace{0.5em}\noindent{\bfseries\color{blue!80!black}【#1】}\par\vspace{0.3em}%
    \begingroup\color{blue!80!black}\emergencystretch=2.5em\BODY\endgroup\par\vspace{0.5em}%
  }

  % 教师专用教学要点、易错警示与评分提示环境（带弹性断行）
  \NewEnviron{teacherNote}[1][教学要点与易错分析]{%
    \par\vspace{0.4em}\noindent{\bfseries\color{teal!80!black}【#1】}\par\vspace{0.2em}%
    \begingroup\color{teal!80!black}\small\emergencystretch=2.5em\BODY\endgroup\par\vspace{0.4em}%
  }
\else
  % 学生版彻底剥离答案与教师批注
  \NewEnviron{solution}[1][]{}
  \NewEnviron{teacherNote}[1][]{}
\fi

% =============================================================================
% 评分量表辅助命令
% =============================================================================

\newcommand{\scoringcriterion}[4]{%
  \textbf{#1 (#2\%):} #3%
  \teacheronly{\quad (\textcolor{blue!70!black}{#4})}%
}

\newcommand{\banddescription}[4]{%
  \item \textbf{#1 分：}#2%
  \teacheronly{\quad$\rightarrow$ #3}%
}

% =============================================================================
% 写作检查清单
% =============================================================================

\newcommand{\writingcheckliststart}{%
  \textbf{检查清单：}%
  \begin{itemize}[leftmargin=*, label=$\square$]%
}
\newcommand{\writingchecklistend}{\end{itemize}}

% =============================================================================
% 表格支持
% =============================================================================

% 特殊表格样式
\newcolumntype{L}[1]{>{\raggedright\arraybackslash}p{#1}}
\newcolumntype{C}[1]{>{\centering\arraybackslash}p{#1}}
\newcolumntype{R}[1]{>{\raggedleft\arraybackslash}p{#1}}
